\section{原始 mirror 子问题的 Gibbs 解}

令
\begin{equation*}
 Z_k(s)=\sum_{a\in\mathcal V} q_k(a\mid s)\exp(\eta r_k(s,a)).
\end{equation*}
由于 $r_k=\log t-\log q_k$，有
\begin{equation}\label{eq:token-partition}
 Z_k(s)=\sum_a q_k(a\mid s)^{1-\eta} t(a\mid s)^\eta.
\end{equation}

\begin{theorem}[Token mirror target]
\label{thm:token-target}
若有限词表上 $ q_k(a\mid s)>0$ 且 $ t(a\mid s)>0$ 对每个 action 成立，则式~\eqref{eq:common-objective} 的唯一最大点为
\begin{equation}\label{eq:token-target}
  p_k^{\mathrm G}(a\mid s)=
 \frac{ q_k(a\mid s)\exp(\eta r_k(s,a))}{Z_k(s)}
 =\frac{ q_k(a\mid s)^{1-\eta} t(a\mid s)^\eta}{Z_k(s)}.
\end{equation}
并且对任意 $p$，
\begin{equation}\label{eq:variational-identity}
 \mathcal J_{k,s}(p)
 =\frac{1}{\eta}\log Z_k(s)
 -\frac{1}{\eta}\KL(p\Vert  p_k^{\mathrm G}).
\end{equation}
\end{theorem}

\begin{proof}
由式~\eqref{eq:token-target}，
\begin{equation*}
 \log  p_k^{\mathrm G}(a\mid s)=\log q_k(a\mid s)+\eta r_k(s,a)-\log Z_k(s).
\end{equation*}
因此
\begin{align*}
 \KL(p\Vert  p_k^{\mathrm G})
 &=\sum_a p(a\mid s)\left[
 \log\frac{p(a\mid s)}{ q_k(a\mid s)}
 -\eta r_k(s,a)+\log Z_k(s)\right]\\
 &=\KL(p\Vert q_k)-\eta\E_p[r_k]+\log Z_k(s).
\end{align*}
移项得到式~\eqref{eq:variational-identity}。KL 非负且仅在 $p= p_k^{\mathrm G}$ 时为零，所以最大点唯一。\hfill$\square$
\end{proof}

\begin{corollary}[几何插值与极限]
当 $\eta\to0^+$ 时 $ p_k^{\mathrm G}\to q_k$；当 $\eta=1$ 时 $ p_k^{\mathrm G}= t$；当 $0<\eta<1$ 时，logits 满足
\begin{equation*}
 \log  p_k^{\mathrm G}=(1-\eta)\log q_k+\eta\log t-\log Z_k.
\end{equation*}
因此 $\eta$ 是 token-level teacher interpolation，而不是 inner optimizer 的学习率。
\end{corollary}

\section{Reward 中心化与回归残差}

定义 $\bar r_k(s)=\mathbb E_{a\sim q_k}[r_k(s,a)]=-\mathcal K_k^T(s)$。将 reward 换成 $r_k-\bar r_k$ 后，配分函数为
\begin{equation*}
\mathbb E_{ q_k}[e^{\eta(r_k-\bar r_k)}]=e^{-\eta\bar r_k}Z_k.
\end{equation*}
分子和配分函数同乘一个 action 无关的常数，原始 mirror 问题的 Gibbs 解不变。下面的平方回归需要区分完整配分函数与均值近似。

三个 residual 全部采用 student 减 target 的方向：
\begin{align}
R^{\mathrm{Part}}_{k,n}(s,a)&\coloneqq U_{k,n}(s,a)-\eta r_k(s,a)+\log Z_k(s),\label{eq:part-residual}\\
R^{\mathrm{Mean}}_{k,n}(s,a)&\coloneqq U_{k,n}(s,a)-G_k(s,a),\label{eq:mean-residual}\\
R^{\mathrm{BPO}}_{k,n}(s,a)&\coloneqq U_{k,n}(s,a)-G_k(s,a)+\mathcal K_{k,n}^S(s).
\end{align}
PMD-Part 的零残差对应 $ p_{k,n}= p_k^{\mathrm G}$。PMD-Mean 以 $\eta\bar r_k$ 近似 $\log Z_k$，仅中心化 reward。对非恒定 reward，由严格 Jensen 不等式有 $\mathbb E_{ q_k}[e^{G_k}]>1$，所以拟合值 $U=G_k$ 不满足概率归一化，最优残差不能为零。其总体最优分布记为 $ p_k^{\mathrm M}$，由第二章的 Lambert-$W$ 解确定。

BPO 则中心化完整的 student--target residual。其 residual 均值为零，而 Mean 的 residual 均值为 $-\mathcal K_{k,n}^S$。这些定义直接区分两个目标，无需额外改变 residual 符号。

\section{参数化策略的投影}

实际神经网络只能表示参数化策略族。固定 prefix 测度 $d_k$ 后，PMD-Part 的总体 inner loss 定义为
\begin{equation}\label{eq:part-loss}
 \mathcal L_{\mathrm{Part}}(\theta_{k,n})
 =\frac12\E_{s\sim d_k,a\sim q_k}
 \left[R^{\mathrm{Part}}_{k,n}(s,a)^2\right].
\end{equation}
若模型在每个 history 上可以实现式~\eqref{eq:token-target}，则全局最小值为零；否则最小值是 anchor history 测度下的函数类投影误差。inner steps 增加只能逼近该投影，不能消除模型表达误差或 history coverage 误差。
